% Source: https://artofproblemsolving.com/community/c3h1365644_right_angle_triangle
% Stable identity: aops-topic-1365644-q01
Let $\angle A=\alpha$ and $AC=b$. In the right triangles $ABC$ and $ABH$, $AB=b\cos\alpha$ and $AH=AB\cos\alpha=b\cos^2\alpha$; similarly, since $\angle C=90^\circ-\alpha$, $BC=b\sin\alpha$ and $CH=BC\sin\alpha=b\sin^2\alpha$.

\textbf{(a)} $AH=CH+AB$ would mean $\cos^2\alpha=\sin^2\alpha+\cos\alpha$, that is $\cos2\alpha=\cos\alpha$, or $2\cos^2\alpha-\cos\alpha-1=0$, so $(2\cos\alpha+1)(\cos\alpha-1)=0$. Then $\cos\alpha=1$ or $\cos\alpha=-\dfrac{1}{2}$, and neither is possible for an acute angle $0^\circ<\alpha<90^\circ$. So no triangle satisfies the condition as posted.

\textbf{(b)} $CH=AH+AB$ means $\sin^2\alpha=\cos^2\alpha+\cos\alpha$, so $-\cos2\alpha=\cos\alpha$, i.e. $2\cos^2\alpha+\cos\alpha-1=0$, $(2\cos\alpha-1)(\cos\alpha+1)=0$. As $\cos\alpha>0$, $\cos\alpha=\dfrac{1}{2}$ and $\alpha=60^\circ$.

\begin{asy}
usepackage("amsmath");
size(5.5cm);
defaultpen(fontsize(10pt));
path ra(pair P,pair u,pair v,real s=0.2){return P+s*unit(u)--P+s*unit(u)+s*unit(v)--P+s*unit(v);}
pair B=(0,0), A=(0,2), C=(2*sqrt(3),0);
pair H=A+dot(B-A,unit(C-A))*unit(C-A);
draw(A--B--C--cycle,linewidth(0.9));
draw(B--H,red+dashed);
draw(ra(B,A-B,C-B,0.18)); draw(ra(H,A-H,B-H,0.15));
draw(arc(A,0.35,-90,-30)); label("$60^\circ$",A+0.55*dir(-60),fontsize(8pt));
draw(arc(C,0.6,150,180)); label("$30^\circ$",C+0.85*dir(165),fontsize(8pt));
dot(H);
label("$A$",A,N); label("$B$",B,SW); label("$C$",C,SE); label("$H$",H,NE);
label("\mbox{}",(min(currentpicture,user=true)+max(currentpicture,user=true))/2,black); // keeps the text after the figure black when LaTeX typesets the labels inline
\end{asy}

So $\boxed{\angle A=60^\circ,\ \angle B=90^\circ,\ \angle C=30^\circ}$. Check with $AC=4$: $AB=2$, $AH=1$, $CH=3=1+2$.
